\documentclass[11pt]{article}
\usepackage[margin=0.9in]{geometry}
\usepackage{amsmath,amssymb,amsthm,booktabs}
\usepackage[T1]{fontenc}
\usepackage{lmodern,microtype}
\usepackage[colorlinks=true,linkcolor=blue,urlcolor=blue]{hyperref}
\newtheorem{theorem}{Theorem}[section]
\newtheorem{lemma}[theorem]{Lemma}
\newtheorem{proposition}[theorem]{Proposition}
\newtheorem{corollary}[theorem]{Corollary}
\theoremstyle{remark}
\newtheorem{remark}[theorem]{Remark}
\newcommand{\R}{\mathbb R}
\newcommand{\C}{\mathbb C}
\newcommand{\norm}[1]{\left\|#1\right\|}
\newcommand{\diag}{\operatorname{diag}}
\newcommand{\ran}{\operatorname{ran}}
\newcommand{\HH}{\mathcal H}
\title{Toward a four-dimensional Yang--Mills construction\\
\large Notebook 12: fine regularity on the flat current domain}
\author{Research working document}
\date{16 October 2026}
\begin{document}
\maketitle
\begin{abstract}
We strengthen the flat fine inverse from physical-cell Hilbert bounds
to pointwise regularity for arbitrary bounded fine forcing. A discrete
heat-kernel argument and the canceled finite-rank formula give uniform
$C^{1,\alpha}$ estimates for every $\alpha<1$, exponential cell locality,
and a componentwise Laplacian bound at zero current. An explicit
counterexample rules out the endpoint and individual second-difference
bounds. Scalar gauge identities then place the full current
perturbation in the actual homogeneous source norms $N$ and $N_\beta$,
without differentiating the current. A smaller uniform complex domain
supports the full flat inverse and extension on bounded forcing.
A massive heat-kernel subtraction also proves right-divergence
estimates on H\"older data. Splitting the time integral at the
output separation removes the exponent loss for each fixed
exponent strictly between zero and one. A separate scalar proof
applies under an explicit positive-penalty hypothesis.
The distinct constrained-current inverse has
the same bounded-forcing regularity on a further common ball.
Nonflat backgrounds and retained histories remain separate.
\end{abstract}

\noindent\textbf{Status.}
Research proofs checked by separate internal derivations and audits.
No external certification or novelty claim is made. The baseline is
valid for every odd $M$. The nonlinear constraint-current term uses
$M=L^j$ with fixed odd $L\ge3$, as proved in Notebook 11. All norms
use physical fine derivatives and volume, at final block scale one.

\section{Previously established algebraic inputs}

We use the exact homogeneous inverse and scalar source definitions of
Notebook 09, and the corrected quadratic-composition bound of Notebook 11.
Let $D$ be the fine scalar gradient, $\mathcal L=D^*D$, $Q_0$ the
normalized scalar block average, and $Q=Q_M^{\rm st}$ the normalized
straight one-form average. The scalar gauge is
\[
 G'=(\mathcal L+Q_0^*a'Q_0)^{-1},\qquad
 R=I-G'Q_0^*(Q_0G'^2Q_0^*)^{-1}Q_0G',\qquad Z=G'RG'.
\]
The distinct positive penalty $a$ enters the fine one-form inverse.
The scalar penalty $a'$ does not replace it. The identities
\[
 Q_0Z=0,\quad \mathcal L G'R=R,\quad
 Z\mathcal L=G'R,\quad \mathcal L Z\mathcal L=R
\]
were proved from these definitions. They will also justify the
fixed-refinement operator products on bounded forcing.

The canceled Bloch formulas $G_0=F+\beta E_a^{-1}\beta^-$,
including every central entry, are those proved in Notebook 09.
We reuse their positive common strip and uniform coefficient bounds;
no new pointwise assertion is assumed from their Hilbert estimates.
For the scalar notation below, with $p_r=z+2\pi r$, set
\[
 a_r=\prod_\nu M^{-1}\sum_{s\in I_M}e^{ip_{r,\nu}s/M},\qquad
 \omega_{r,\mu}=a_rM^{-1}\sum_{t=0}^{M-1}e^{ip_{r,\mu}t/M},
\]
\[
 R_{s,2}=\sum_{r\ne0}\frac{a_ra_r^-}{P_r^2},\qquad
 D_{s,2}=a_0a_0^-+P_0^2R_{s,2}.
\]
The symbols $P_r$ and the negative-argument convention are recalled
in the next section. These are scalar denominators; the subscript
$s,2$ is a name and does not designate a free matrix index.
The homogeneous source value of the one-form penalty is $a=1$;
a fixed compact positive interval is allowed throughout.
% Pointwise flat inverse; no standalone preamble.
% Requires amsmath/amssymb/amsthm, theorem environment, and \norm.
% Every label introduced by this fragment has the unique prefix pt-.
\section{Pointwise regularity for arbitrary fine forcing}
\label{pt-section}

The full flat inverse also obeys pointwise estimates. These estimates
concern arbitrary bounded fine forcing, without assuming that it is a
coarse lift. Retain fine spacing $\eta=M^{-1}$, odd $M$, unit coarse
cells, the physical fine measure $\eta^d$, and
\begin{equation}\label{pt-operator}
 \mathcal A_a=d^*d+DRD^*+Q^*aQ,\qquad G_0=\mathcal A_a^{-1},
 \qquad \mathcal L=D^*D.
\end{equation}
Here $d$ is the fine curl, $D$ the scalar gradient, $Q=Q_M^{\rm st}$,
and $R$ is the matched homogeneous flat scalar source projection.
The penalty lies in a fixed interval $0<a_-\le a\le a_+<\infty$;
the normalized source value $a=1$ is included. The discussion applies
to the homogeneous infinite lattice and compatible periodic regulators.

For a fixed coarse-cell patch $U$, define the discrete pointwise norm
\begin{equation}\label{pt-holder-norm}
 \norm{u}_{1,\alpha;U}=\norm{u}_{\infty;U}
  +\sum_\mu\norm{D_\mu^+u}_{\infty;U}
  +\sum_\mu\sup_{\substack{x,y\in U\\x\ne y}}
       \frac{|D_\mu^+u(x)-D_\mu^+u(y)|}{|x-y|_1^\alpha}.
\end{equation}
The points are actual fine-grid points and a fixed collar supplies the
differences. Enlarging the patch by finitely many neighboring coarse
cells changes only constants. Backward differences obey the same bounds.

\begin{theorem}\label{pt-theorem}
For each $0<\alpha<1$ there are $C_\alpha,\rho>0$, independent of
odd $M$ and compatible periodic volume, such that the block from
forcing in coarse cell $m$ to a fixed collar of coarse cell $n$ satisfies
\begin{align}
 \norm{G_0(n,m)}_{\ell^\infty\to C^{1,\alpha}}
       &\le C_\alpha e^{-\rho|n-m|_1},
       \label{pt-pointwise-bound}\\
 \norm{(\mathcal L G_0)(n,m)}_{\ell^\infty\to\ell^\infty}
       &\le C e^{-\rho|n-m|_1}.
       \label{pt-laplacian-bound}
\end{align}
The first target norm is \eqref{pt-holder-norm}; $\mathcal L$ in
the second estimate acts componentwise on the vector output.
On tori use periodic coarse-cell distance. Weighted row and column
block sums are uniformly bounded at every smaller positive exponent.
Consequently bounded fine forcing produces uniformly bounded fields,
first differences, locally $\alpha$-H\"older first differences, and
componentwise Laplacians.
\end{theorem}

The constants need not stay bounded as $\alpha\uparrow1$. In fact a
uniform Lipschitz estimate for the first differences, and a uniform
bound for each individual second difference, are false in dimension
$d\ge2$; an example is given below. No source weak-norm identification
is inferred merely from the theorem's terminology.

\subsection{Reduction to a scalar high-mode inverse}

Use the aliases $p_r=z+2\pi r$, $r\in I_M^d$, with
$I_M=\{-(M-1)/2,\ldots,(M-1)/2\}$, and write
\begin{equation}\label{pt-alias-data}
 u_{r,\mu}^{\pm}=M(e^{\pm ip_{r,\mu}/M}-1),\qquad
 P_r=\sum_\mu u_{r,\mu}^+u_{r,\mu}^-,\qquad
 \mathcal F_M(z)_{sr}=e^{i(z+2\pi r)\cdot s/M}.
\end{equation}
The common holomorphic chart from the full inverse theorem contains
the real fundamental cube in the interior of $|\operatorname{Re}z_\mu|<K<2\pi$.
For real $z$, $\mathcal F_M$ is unitary from alias counting $\ell^2$
to physical-cell $L^2$. On the complex chart it and its inverse have
uniformly bounded $L^2$ norms. In particular a bounded fine-cell input
has alias coefficients bounded in $\ell^2$, and individually, by
$C\norm{f}_\infty$.

Recall the canceled inverse $G_0=F+\beta E_a^{-1}\beta^-$.
Separate the scalar orientation-diagonal multiplier
\begin{equation}\label{pt-high-scalar}
 T_r=P_r^{-1}\quad(r\ne0),\qquad T_0=0.
\end{equation}
The difference $F-T$ is a fixed sum of rank-one operators per vector
orientation: its central block, two mixed blocks, and one high/high
term. Its high vectors are $\omega_{r,\mu}^-/P_r$ and the analytic
transpose partners. The remaining scalar coefficients are uniformly
bounded and holomorphic. Thus the only full-rank pointwise problem is
\eqref{pt-high-scalar}. Minus signs on multipliers throughout this
argument denote evaluation at negative full momentum, with the
analytic transpose; they do not conjugate a complex variable.

\subsection{Weighted heat estimates with complex Bloch twist}

Let $p_s(n)$ be the continuous-time nearest-neighbor heat kernel on
$\mathbb Z$ with generator $f(n+1)+f(n-1)-2f(n)$. Its Poisson-process
representation gives
\begin{equation}\label{pt-heat-moment}
 \sum_n e^{tn}p_s(n)=e^{2s(\cosh t-1)},\qquad
 \sum_n p_s(n)=1.
\end{equation}
It is symmetric and decreasing for $n\ge0$. For completeness, set
$q_n=p_s(n)-p_s(n+1)$ on $n\ge0$. Then
\[
 q_0'=q_1-3q_0,\qquad
 q_n'=q_{n-1}+q_{n+1}-2q_n\quad(n\ge1),
\]
with nonnegative initial data. This bounded cooperative system
preserves nonnegativity, as follows by shifting its generator to a
matrix with nonnegative entries and expanding the exponential.
Also Fourier inversion gives
\begin{equation}\label{pt-heat-origin}
 p_s(0)=\frac1{2\pi}\int_{-\pi}^{\pi}
               e^{-4s\sin^2(t/2)}\,dt
       \le C(1+s)^{-1/2}.
\end{equation}
For $0\le r\le r_0$, summation by parts therefore yields the exact
weighted variation identity
\begin{align}
 \sum_n e^{r|n|}|p_s(n+1)-p_s(n)|
  &=(1+e^r)\left[p_s(0)
       +(1-e^{-r})\sum_{n\ge1}e^{rn}p_s(n)\right]\notag\\
  &\le C\left[(1+s)^{-1/2}
                   +r e^{2s(\cosh r-1)}\right].
       \label{pt-weighted-variation}
\end{align}
Take $s=M^2t$ and $r=\delta/M$, where $\delta$ is a fixed small
constant. Multiplying the difference by $M$ gives
\begin{equation}\label{pt-first-heat}
 \sum_n e^{\delta|n|/M}
       |D_\eta p_{M^2t}(n)|
       \le C(t^{-1/2}+\delta)e^{C\delta^2t}.
\end{equation}
For $t<\eta^2$ the trivial $M$ bound is smaller than $t^{-1/2}$.
The weight is submultiplicative. Apply
\eqref{pt-first-heat} twice to
$p_s=p_{s/2}*p_{s/2}$ to obtain the second-difference bound.
Tensor products then give, for any specified composition of $k$ fine
forward/backward differences in $d$ dimensions,
\begin{equation}\label{pt-heat-two-derivatives}
 \sum_{n\in\mathbb Z^d}e^{\delta|n|_1/M}
       |D_\eta^k p_{M^2t}(n)|
       \le Ct^{-k/2},\qquad 0<t\le1,\quad k=0,1,2.
\end{equation}
The displayed $p$ is a probability kernel. If it is instead written
as a density against physical fine volume, the density factor $M^d$
cancels the summation factor $M^{-d}$.

The reconstructed Bloch heat operator has the exact image-sum kernel
\begin{equation}\label{pt-twisted-heat}
 H_z(t;s,s')=\sum_{k\in\mathbb Z^d}
             e^{-iz\cdot k}p_{M^2t}(s-s'+kM).
\end{equation}
This convention uses ordinary summation over the fine-cell input
index $s'$. The formula follows by finite Fourier inversion and
continues to complex $z$ by the exponential moment bound.
For $s$ in a fixed output collar and $s'$ in the input cell,
\[
 |e^{-iz\cdot k}|
       \le C e^{\delta|s-s'+kM|_1/M}.
\]
Thus \eqref{pt-heat-two-derivatives}, including the corresponding
output differences of the image sum, proves
\begin{equation}\label{pt-twisted-derivatives}
 \norm{D_\eta^kH_z(t)}_{\infty\to\infty}
       \le Ct^{-k/2},\qquad 0<t\le1,\quad k=0,1,2.
\end{equation}
The difference between first-gradient rows at fine points $x,y$ is
bounded both by $Ct^{-1/2}$ and, by telescoping, by
$C|x-y|_1t^{-1}$. Interpolating these two elementary bounds gives
\begin{equation}\label{pt-heat-holder}
 \norm{H_z(t)}_{\infty\to C^{1,\alpha}}
       \le C_\alpha t^{-(1+\alpha)/2},\qquad 0<t\le1.
\end{equation}
Subtracting the central heat mode changes this only by a bounded
term, since $P_0,u_0$ and its plane-wave profile are uniformly bounded
on the chart. For $t\ge1$, the high-alias bound
$\operatorname{Re}P_r\ge c|r|^2$ and the individual input-coefficient
bound instead give
\begin{align}
 &\norm{H_z(t)-e^{-tP_0}\Pi_0(z)}_{\infty\to C^{1,\alpha}}
       \notag\\
 &\qquad\le C\sum_{r\ne0}(1+|r|)^{1+\alpha}e^{-ct|r|^2}
       \le C_\alpha e^{-ct/2}.
       \label{pt-heat-large-time}
\end{align}
The high-mode inverse is therefore the convergent operator integral
\begin{equation}\label{pt-scalar-integral}
 T(z)=\int_0^\infty
             \bigl[H_z(t)-e^{-tP_0}\Pi_0(z)\bigr],dt.
\end{equation}
The small-time integral is finite because
$\int_0^1t^{-(1+\alpha)/2}\,dt=2/(1-\alpha)$.
This proves a uniform holomorphic
$\ell^\infty\to C^{1,\alpha}$ bound for $T$ for every $\alpha<1$.

\subsection{Pointwise bounds for the rank corrections}

The centered scalar and straight-vector averaging multipliers satisfy
\begin{equation}\label{pt-alias-weights}
 |a_r^\pm|\le Cw(r),\qquad
 |\omega_{r,\mu}^\pm|\le\frac{Cw(r)}{1+|r_\mu|},\qquad
 w(r)=\prod_\nu(1+|r_\nu|)^{-1}.
\end{equation}
Indeed, for a nonzero component, the numerator of
$\sin((z_\nu+2\pi r_\nu)/2)/
 [M\sin((z_\nu+2\pi r_\nu)/(2M))]$
is uniformly bounded and its denominator has modulus at least
$c|r_\nu|$ on the enlarged chart. The zero component is bounded
directly. The vector multiplier has one extra longitudinal factor.
This includes aliases at the edge of the fine zone.

On a dyadic maximum-norm shell,
\begin{equation}\label{pt-shell-sum}
 \sum_{2^j\le|r|_\infty<2^{j+1}}w(r)
                      \le C_d(j+1)^{d-1}.
\end{equation}
Choose a coordinate at the maximal scale. Its reciprocal weight sums
to $O(1)$ over that shell, and each other coordinate contributes
$O(j+1)$. Summing over the possible maximal coordinates proves
\eqref{pt-shell-sum}. Consequently
\begin{equation}\label{pt-rank-holder-sum}
 \sum_{r\ne0}(1+|r|)^{1+\alpha}
             \left|\frac{\omega_{r,\mu}^-}{P_r}\right|
   \le C\sum_{r\ne0}w(r)(1+|r|)^{\alpha-1}
   \le C_\alpha.
\end{equation}
Absolute Fourier summation now bounds the reconstructed high rank
vectors in $C^{1,\alpha}$. Their central partners have the same bound.
The Woodbury vector has
\begin{equation}\label{pt-beta-bound}
 |\beta_{r,\mu}|
  \le C\left(\frac{|a_r^-|}{(1+|r|)^3}
                  +\frac{|\omega_{r,\mu}^-|}{(1+|r|)^2}\right),
       \qquad r\ne0,
\end{equation}
so its output profile satisfies the same weighted sum. Its canceled
central component is bounded directly.

The input rows of every rank term have uniformly bounded alias
$\ell^2$ norm. By the physical-cell Fourier normalization,
\[
 |\text{input row}(f)|
       \le C\norm{f}_{L^2(\text{physical cell})}
       \le C\norm{f}_\infty.
\]
The cell has physical volume one; there is no factor $M^{d/2}$.
Combining these row estimates with
\eqref{pt-rank-holder-sum}--\eqref{pt-beta-bound} controls all rank
corrections. The full reconstructed $G_0$ therefore has the required
uniform holomorphic pointwise bound on the common chart.

\subsection{The componentwise Laplacian}

For the scalar term, $\mathcal L T=I-\Pi_0$ is bounded on
$\ell^\infty$. Applying $\mathcal L$ to a high rank vector
$\omega_{r,\mu}^-/P_r$ gives $\omega_{r,\mu}^-$. Its full
reconstruction is a column of the physical adjoint average, minus
its bounded central mode. Explicitly, for $x=n+s/M$, $M=2h+1$,
\begin{equation}\label{pt-adjoint-average}
 (Q^*B)_\mu(x)=b_s B_\mu(n)+(1-b_s)B_\mu(n-e_\mu),
 \qquad b_s=\frac{s_\mu+h+1}{M}.
\end{equation}
These weights are nonnegative and sum to one. The Bloch profile is
$b_s+(1-b_s)e^{-iz_\mu}$, uniformly bounded on the complex chart
and any fixed collar. An absolute alias sum of $\omega$ is unnecessary
and would introduce a spurious logarithmic loss.

For the Woodbury vector, its high-alias formula gives
\begin{equation}\label{pt-beta-laplacian}
 (\mathcal L\beta)_{r,\mu}
  =\frac{P_0u_{r,\mu}^+a_r^-}{P_rD_{s,2}}
           -\omega_{r,\mu}^-k_\mu d_\mu^+.
\end{equation}
The first term is $O(w(r)/(1+|r|))$, absolutely summable by
\eqref{pt-shell-sum}. The second is a bounded coefficient times
the profile \eqref{pt-adjoint-average}, with its central subtraction.
The canceled central Laplacian component is bounded. The same input
row estimates as before now prove a uniform holomorphic
$\ell^\infty\to\ell^\infty$ bound for $\mathcal L G_0$.
This argument supplies neither H\"older continuity of its Laplacian
nor bounds on all separate second differences. In particular the
$Q^*$ profile can have transverse block jumps.

\subsection{Gluing, exponential decay, and bounded forcing}

The whole reconstructed inverse and its fine derivatives are periodic
in coarse momentum. The full real symbol is invariant under the
alias relabeling induced by $z\mapsto z+2\pi n$; the within-cell
integer phase cancels wrapping of the alias index. The enlarged
holomorphic charts agree on their real overlaps and hence glue.
Individual separated central/high terms need not be periodic; they
were used only to bound the whole symbol on each chart.

Output collars consist of finitely many neighboring-cell copies,
with bounded holomorphic factors $e^{iz\cdot n}$. Thus the common
strip bounds hold as operator bounds into the discrete norm
\eqref{pt-holder-norm}. For each fixed $M$ these are finite-dimensional
Banach spaces, so entrywise holomorphy and the proved uniform norms
give Banach-valued contour shifting. This proves
\eqref{pt-pointwise-bound}--\eqref{pt-laplacian-bound} without a
growing-dimension factor.

In infinite volume the resulting absolutely summable blocks define
the inverse on bounded forcing and agree with the Hilbert inverse
on finitely supported fine forcing. The inverse identity extends
through absolute convergence of the exponentially summable operator
products at each fixed $M$. The source projection is not local, but
$R=\mathcal L Z\mathcal L$ has exponentially summable cell $L^2$
blocks. Finite-cell norm equivalence makes their pointwise blocks
summable for each fixed $M$. That last identity argument need not
be uniform in $M$: the pointwise inverse bound was proved uniformly
by the heat and rank estimates above.

Finite compatible tori sample the same symbol, giving exact
periodizations and the same bounds in periodic distance, including
short periods and multiplicities. Summing the cell kernels bounds
the field and gradient globally and the gradient H\"older seminorm
on every fixed collar. Nearby points lie in one enlarged collar;
farther points use twice the gradient supremum. This also gives the
global discrete H\"older bound, with periodic physical distance on
tori.

\subsection{A sharp obstruction to the Lipschitz endpoint}

For $d\ge2$, take the compatible torus consisting of one coarse
cell. At its sole coarse momentum, zero, the full inverse is
$a^{-1}$ on the central constant mode and $P_r^{-1}$ on every high
alias. Thus on fine mean-zero forcing it is exactly the scalar
periodic inverse Laplacian, independently in each vector component.

Choose a smooth cutoff $\chi$ supported in a small ball inside the
unit continuum torus and equal to one near its center. Set
\begin{equation}\label{pt-endpoint-example}
 u_\varepsilon(x)=\chi(x)x_1x_2\,
                       \frac12\log(|x|^2+\varepsilon^2).
\end{equation}
For $P=x_1x_2$ and $f(r)=\frac12\log(r^2+\varepsilon^2)$,
harmonicity and homogeneity of $P$ give
\begin{equation}\label{pt-endpoint-laplacian}
 \Delta(Pf)=P\left(f''+\frac{d+3}{r}f'\right)
   =P\frac{(d+2)r^2+(d+4)\varepsilon^2}
                  {(r^2+\varepsilon^2)^2}.
\end{equation}
Since $|P|\le r^2/2$, this is bounded independently of $\varepsilon$.
The cutoff-region derivatives are also uniformly bounded because
that region stays away from zero. However
\begin{equation}\label{pt-endpoint-divergence}
 \partial_1\partial_2u_\varepsilon(0)=\log\varepsilon.
\end{equation}

For fixed $\varepsilon>0$, sample the smooth periodic function on
increasingly fine odd grids and subtract its exact discrete mean.
Define $f_{M,\varepsilon}=\mathcal L_\eta u_{M,\varepsilon}$.
Its mean is zero, and Taylor convergence gives
$\norm{f_{M,\varepsilon}}_\infty\le C+1$ for sufficiently large
$M$, depending on $\varepsilon$. The full inverse returns exactly
$u_{M,\varepsilon}$. Its mixed second fine difference at the origin
converges to \eqref{pt-endpoint-divergence}. Choosing small
$\varepsilon$ and then large $M$ proves that bounded forcing can
produce arbitrarily large mixed second differences. A uniform
Lipschitz bound on first differences would bound these differences
and is therefore impossible. This does not contradict
\eqref{pt-laplacian-bound}, which controls their Laplacian sum rather
than each second difference separately.

\subsection{Source scope}

The theorem supplies arbitrary-forcing pointwise estimates for the
matched homogeneous flat inverse, including its normalized $a=1$
case. It does not by itself identify every numerical factor or term
in the source's weak norms, establish nonflat estimates, handle
general retained-history domains, or control the complete interacting
current carrier. Those applications require exact norm and coordinate
checks. The false endpoint above must not enter such a match through
an implicit assertion that $C^{1,\alpha}$ controls all second derivatives.
% Pointwise full-current neighborhood; no standalone preamble.
% Requires amsmath/amssymb/amsthm, theorem environment, and \norm.
% Uses the preceding pointwise baseline theorem. Labels have prefix pc-.
\section{The full flat current in the pointwise weak norms}
\label{pc-section}

We now transfer the pointwise estimates to arbitrary bounded complex
fine currents. The background remains exactly flat. For the corrected
nonlinear constraint term, refinement is $M=L^j$ with fixed odd $L$,
as in the independently audited fixed-step quadratic-composition
theorem. The result does not use a pointwise bound for the projection
$\Pi$ on arbitrary bounded fields and does not differentiate the current.

\subsection{The source norms and spatial weights}

At homogeneous $\ell=1$ and background $U=1$, deposited Raw I,
lines 12740--12752, defines the equivalent norms
\begin{equation}\label{pc-source-norms}
 N(A)=\max\{\norm{A}_\infty,\norm{\nabla A}_\infty\},\qquad
 N_\beta(A)=\max\{N(A),[\nabla A]_\beta\},
 \quad 0\le\beta\le\beta_0<1.
\end{equation}
Here $\nabla$ collects the first fine differences of a vector field;
the H\"older member can be taken on fixed coarse-cell collars.
The pointwise estimates of the preceding section also imply its
global discrete version. There is no Laplacian in $N_\beta$.
The source explicitly excludes the Laplacian member for the current
inverse at Raw I, lines 13759--13765.

Let $\varpi(n)>0$ be a cell weight satisfying
\begin{equation}\label{pc-weight-ratio}
 \frac{\varpi(n)}{\varpi(m)}
       \le e^{\kappa\operatorname{dist}(n,m)}.
\end{equation}
The distance is ordinary or periodic cell distance. Define
$\norm{f}_{\infty,\varpi}$, $N_\varpi$, and $N_{\beta,\varpi}$
by the corresponding weighted supremum over cells and fixed collars.
Neighboring weights have uniformly bounded ratios. In particular
$\varpi(n)=e^{\kappa\operatorname{dist}(n,m_0)}$ is allowed.
These norms are applied to complete fields. We never sharply cut a
differentiable field into cells before applying $D^*$, which would
introduce artificial fine-scale boundary derivatives.

Retain
\begin{equation}\label{pc-scalar-data}
 \mathcal L=D^*D,\qquad Z=G'RG',\qquad
 \Lambda=Z\mathcal L D^*=G'RD^*,\qquad \Pi=I-D\Lambda.
\end{equation}
Stars denote physical adjoints on the real slice and their fixed
analytic bilinear continuations for complex parameters. The
penalized baseline is
$G_0=(d^*d+DRD^*+Q^*aQ)^{-1}$, with matched unit-scale value $a=1$.

\subsection{The scalar smoothing needed for the current}

For a common small $\kappa>0$, the scalar inverse supplies constants
independent of $M$ and compatible volume such that
\begin{align}
 \norm{Z\mathcal L f}_{\infty,\varpi}
       &\le c_0\norm{f}_{\infty,\varpi},&
 \norm{DZ\mathcal L f}_{\infty,\varpi}
       &\le c_1\norm{f}_{\infty,\varpi},
       \label{pc-scalar-bounded}\\
 N_{\beta,\varpi}(DZf)
       &\le z_1\norm{f}_{\infty,\varpi},&
 N_{\beta,\varpi}(DZD^*f)
       &\le z_2\norm{f}_{\infty,\varpi}.
       \label{pc-scalar-holder}
\end{align}
Constants can be chosen at the fixed ceiling $\beta_0$, uniformly
for $0\le\beta\le\beta_0$.

Here is the pointwise proof of these additional scalar facts. The
canceled alias matrix of $Z$ is the high diagonal $P_r^{-2}$ plus
a fixed number of rank terms whose high vectors are $a_r^-/P_r^2$
and their transpose partners. Its central entries have no massless
pole. The high diagonal has the heat representation
\begin{equation}\label{pc-biharmonic-heat}
 \int_0^\infty t\,[H_z(t)-e^{-tP_0}\Pi_0(z)]\,dt.
\end{equation}
The weighted heat derivative bounds extend to every fixed integer
order $k$: split the semigroup into $k$ equal-time factors, put one
fine difference on each, and use weighted $L^1$ convolution.
Thus the small-time bound is $C_kt^{-k/2}$. An additional
$\beta$-H\"older output difference costs $t^{-\beta/2}$ by the
same interpolation used for the baseline inverse. Input derivatives
are derivatives of the kernel, or discrete summation by parts, and
do not require derivatives of the forcing. The relevant integral is
\[
 \int_0^1 t^{1-(m+\beta)/2}\,dt<\infty
                  \quad\text{if }m+\beta<4.
\]
For $Z\mathcal L,DZ\mathcal L$ the derivative orders are two and
three. For the gradient H\"older members of $DZ,DZD^*$ they are
two and three, respectively, before the extra $\beta$. Large times
use the high-alias gap.

For the rank terms, an input row with the required derivatives has
uniform alias $\ell^2$ norm. Output vectors obey
$|a_r/P_r^2|\le Cw(r)/(1+|r|)^4$, where
$w(r)=\prod_\nu(1+|r_\nu|)^{-1}$ and its dyadic shell sum is
$O((j+1)^{d-1})$. Every output derivative/H\"older order strictly
below four is therefore absolutely summable. Physical-cell $L^2$
input rows act on bounded forcing without a dimension factor.
The same periodic holomorphic gluing as for $G_0$ gives exponential
cell kernels and hence \eqref{pc-scalar-bounded}--\eqref{pc-scalar-holder}.
Finite periodic sampling preserves the bounds.

It follows, using the actual derivative of the complete field, that
\begin{equation}\label{pc-lambda-on-N}
 \norm{\Lambda A}_{\infty,\varpi}\le l_0N_\varpi(A),\qquad
 \norm{D\Lambda A}_{\infty,\varpi}\le l_1N_\varpi(A),\qquad
 \norm{\Pi A}_{\infty,\varpi}\le(1+l_1)N_\varpi(A).
\end{equation}
No assertion that $D\Lambda$ or $\Pi$ is bounded on unrestricted
$\ell^\infty$ input is used.

\subsection{Both current terms map $N$ to $N_\beta$}

The matched local auxiliary factorization is
\begin{equation}\label{pc-aux-factorization}
 \Theta'_J=\mathfrak j_J^*\Lambda+\Lambda^*\mathfrak j_J
       +\Lambda^*V_JD\Lambda+(D\Lambda)^*V_J^*\Lambda.
\end{equation}
The finite-incidence coefficient operators $\mathfrak j_J,V_J$
and their analytic transposes have weighted pointwise norms at most
$c_j\norm{J}_\infty,c_v\norm{J}_\infty$. They contain no derivative
of $J$. The exact scalar identity $G_0\Lambda^*=DZ$ gives
\begin{equation}\label{pc-outer-aux-smoothing}
 G_0\Theta'_J
  =G_0\mathfrak j_J^*\Lambda+DZ\mathfrak j_J
       +DZV_JD\Lambda+DZD^*V_J^*\Lambda.
\end{equation}
Let $g_\beta$ bound $G_0$ from weighted bounded forcing to
$N_{\beta,\varpi}$. Then
\begin{equation}\label{pc-aux-N-bound}
 N_{\beta,\varpi}(G_0\Theta'_JA)
        \le C_{\rm aux}\norm{J}_\infty N_\varpi(A),
\end{equation}
where one possible constant is
\[
 C_{\rm aux}=g_\beta c_jl_0+z_1c_j+z_1c_vl_1+z_2c_vl_0.
\]
This is the homogeneous normalized source (J1), including weights.
Its input requires $N$ only, not H\"older regularity of $A$ or $J$.

For the full constraint term, write
\begin{equation}\label{pc-constraint-weight}
 \Theta_J=T_{v_J},\qquad v_J=E_J^T J,
 \qquad
 \langle A,T_vB\rangle_\eta
      =2\langle C_j^{(2)}(A,B),v\rangle_1,
\end{equation}
where $E_J$ is the auxiliary normalized extension and $C_j^{(2)}$
is the quadratic jet of the actual corrected fixed-$L$ composition.
The previously proved auxiliary cell-Schur estimate gives, analytically,
$\norm{v_J}_\infty\le C_v\norm{J}_\infty$ throughout its current ball.

The direct pointwise estimate for $T_v$ uses the same physical
factor as its Hilbert proof. At a step of fine spacing $h$, the
quadratic coefficient is $(h/L)C_L$. Its finite coefficient tensor,
paired with coarse measure and converted to the fine physical
adjoint, has an absolute row sum at most $c_Lh\norm{v}_\infty$.
The factor $L^d$ in this conversion is fixed. The quadratic chain
rule expresses the full operator as
\begin{equation}\label{pc-pointwise-chain}
 T_v=\sum_{t=0}^{j-1}Q_{0:t}^*T_{h_t,v_t}Q_{0:t},\qquad
 h_t=L^{t-j},\qquad v_t=Q_{t+1:j}^*v.
\end{equation}
Both flat $Q$ and its physical adjoint $Q^*$ have nonnegative row
coefficients summing to one. They are $\ell^\infty$ contractions,
including the covering-space multiplicities on short tori. Thus
\[
 \norm{T_v}_{\infty\to\infty}
       \le\frac{c_L}{L-1}\norm{v}_\infty.
\]
The composition has fixed unit-cell range $R_*$, since its reading
radii sum over the same geometric series. Consequently
\begin{equation}\label{pc-T-pointwise}
 \norm{T_vf}_{\infty,\varpi}
       \le\tau_\infty\norm{v}_\infty
                         \norm{f}_{\infty,\varpi}.
\end{equation}
The constant includes $e^{\kappa R_*}$. This is a coefficient-row
argument, not a conversion of a final-cell Hilbert norm into a
pointwise norm.

The identity
\begin{equation}\label{pc-Pi-outer}
 G_0\Pi^*=G_0-DZD^*
\end{equation}
now places every smoothing operator outside the bounded current
coefficients. Combining \eqref{pc-lambda-on-N},
\eqref{pc-T-pointwise}, and \eqref{pc-Pi-outer} gives
\begin{align}
 N_{\beta,\varpi}(G_0\Pi^*\Theta_J\Pi A)
       &\le C_{\rm con}\norm{J}_\infty N_\varpi(A),
       \label{pc-constraint-N-bound}\\
 C_{\rm con}&=(g_\beta+z_2)\tau_\infty C_v(1+l_1).
       \notag
\end{align}
In particular no derivative of $J$ has appeared.

There is also a constrained baseline estimate. Put
\begin{equation}\label{pc-constrained-baseline}
 \Gamma=G_0-H_{\rm src}q_{\rm aux}^{-1}H_{\rm src}^*,
 \qquad q_{\rm aux}=q_M+aI.
\end{equation}
The extension $H_{\rm src}$ is pointwise local into $N_\beta$;
its transpose maps bounded fine fields to bounded coarse fields
by the physical-cell Schur bound; and
$q_{\rm aux}^{-1}=QG_0Q^*$ is a local bounded coarse operator.
Hence $\Gamma$ maps weighted bounded forcing to $N_\beta$ and
\begin{equation}\label{pc-constrained-current}
 N_{\beta,\varpi}(\Gamma\Theta_JA)
      \le C\norm{J}_\infty\norm{A}_{\infty,\varpi}.
\end{equation}

\subsection{A common pointwise current neighborhood}

\begin{theorem}\label{pc-theorem}
Let $\varepsilon_H$ be the established full-current Hilbert/Schur
radius, which is contained in the auxiliary radius. Define
\begin{equation}\label{pc-radius}
 C_*=C_{\rm aux}+C_{\rm con},\qquad
 \varepsilon_{\rm pt}
       =\min\{\varepsilon_H,(2C_*)^{-1}\},
\end{equation}
increasing $C_*$ by fixed norm-convention factors if necessary.
For $\norm{J}_\infty<\varepsilon_{\rm pt}$, put
\begin{equation}\label{pc-full-perturbation}
 \Psi_J=\Theta'_J+\Pi^*\Theta_J\Pi,\qquad
 K_J=G_0\Psi_J.
\end{equation}
The actual full homogeneous flat inverse is the convergent series
\begin{equation}\label{pc-neumann-inverse}
 G_1^Jf=(I-K_J)^{-1}G_0f
                  =\sum_{n\ge0}K_J^nG_0f.
\end{equation}
It is holomorphic from the complex bounded-current ball to operators
from bounded fine forcing into $N_\beta$, uniformly for
$0\le\beta\le\beta_0$. With $g_N$ the baseline bounded-forcing-to-$N$
constant, it obeys
\begin{align}
 N_\varpi(G_1^Jf)
       &\le2g_N\norm{f}_{\infty,\varpi},
       \label{pc-inverse-N}\\
 N_{\beta,\varpi}(G_1^Jf)
       &\le(g_\beta+2C_*\varepsilon_{\rm pt}g_N)
                              \norm{f}_{\infty,\varpi}.
       \label{pc-inverse-Nbeta}
\end{align}
All constants are independent of $M=L^j$ and compatible volume.
For forcing supported in cell $m$, choosing
$\varpi(n)=e^{\kappa\operatorname{dist}(n,m)}$ gives exponential
coarse-cell locality in the pointwise target norms. The auxiliary
inverse is the same result with only $\Theta'_J$ retained.
\end{theorem}

\begin{proof}
Equations \eqref{pc-aux-N-bound} and
\eqref{pc-constraint-N-bound} give
\[
 N_{\beta,\varpi}(K_JA)
          \le C_*\norm{J}_\infty N_\varpi(A).
\]
The target dominates $N_\varpi$, so the ratio is below $1/2$
on \eqref{pc-radius}. The series converges in weighted $N$.
For every term after the zeroth, apply the stronger $N$-to-$N_\beta$
bound to its outside $K_J$ and the $N$ contraction to the remaining
factors. Summing gives \eqref{pc-inverse-N}--\eqref{pc-inverse-Nbeta}.
The coefficient $\Theta_J=T_{E_J^TJ}$ is nonlinear in $J$ but is
analytic on the auxiliary ball with the proved linear norm bound.
Uniform convergence therefore gives Banach-valued holomorphy; no
linear-current differentiation formula is presumed for the full family.

For forcing in one cell, the exponentially weighted solution is in
physical $L^2$ on the infinite lattice. It consequently equals the
independently constructed Hilbert inverse on that input. Equality
for bounded forcing follows by the absolutely summable pointwise
cell kernels and the fixed-$M$ summability of the original operator
products. The finite-volume proof is identical with periodic
distance. The minimum with $\varepsilon_H$ ensures that both
constructions use the same actual source inverse.
\end{proof}

The normalized extensions inherit the same bounds:
$E_1^J=G_1^JQ^*q_1^J$ and $E_J=G_\pi^JQ^*q_J$.
Here $Q^*$ maps bounded coarse fields to bounded fine fields, and
the established coarse Schur bounds make $q_1^J,q_J$ bounded on
weighted coarse $\ell^\infty$. Thus their output $N_\beta$ norms
are uniform on the common ball.

\subsection{What has and has not matched the source}

The input and output norms in \eqref{pc-aux-N-bound} and
\eqref{pc-constraint-N-bound} are the explicit homogeneous flat
$N,N_\beta$ norms, up to fixed equivalent conventions for localized
H\"older seminorms. The current has only a supremum bound. These
estimates and \eqref{pc-inverse-Nbeta} supply the bounded-forcing
members of the current program and their weighted versions.

No current Laplacian estimate is asserted. The source expressly
excludes that member, and the Laplacian of $DZD^*$ is an order-zero
operator whose bounded-input pointwise norm has not been supplied.
The following sections supply the one-form right-divergence members
using complete weighted H\"older fields and smooth physical
cutoffs, then treat the distinct scalar $G'$ under its stated
penalty hypothesis.

The fixed-step corrected-average source identification remains the
one audited from the deposited transcriptions. These results are
homogeneous and flat. They do not establish nonflat source stability,
mixed retained-domain estimates or an interacting Yang--Mills
ultraviolet limit. Companion Notebook 13 constructs a uniform
analytic chart for the full fixed-reference nonlinear constraint;
its moving-background identification is a further question.

% Right-divergence and constrained-current members; no preamble.
% Requires amsmath/amssymb/amsthm, theorem environment, and \norm.
% Depends on the preceding pt- and pc- results. Labels have prefix rd-.
\section{Right-divergence data and the constrained current inverse}
\label{rd-section}

The remaining right-divergence estimates require cancellation on
H\"older input. We isolate that cancellation in a translation-invariant
massive lattice Green function. The argument applies to the full
homogeneous flat $G_0$ and its auxiliary and full current perturbations.
It does not identify the distinct scalar inverse $G'$ with $G_0$.

Use the admissible spatial weights and weak norms of the preceding
section. Let $\norm{\lambda}_{\gamma,\varpi}$ denote the weighted
supremum of the fine field and its $\gamma$-H\"older seminorm on
fixed cell collars. Fields are kept complete; no H\"older datum is
multiplied by a sharp cell indicator.

\begin{theorem}\label{rd-theorem}
For $0\le\beta<\gamma<1$ and a sufficiently small fixed spatial
exponent $\kappa$, the homogeneous flat inverses satisfy
\begin{equation}\label{rd-main-estimate}
 N_{\beta,\varpi}(G D^*\lambda)
       \le C_{\beta,\gamma}\norm{\lambda}_{\gamma,\varpi},
 \qquad G\in\{G_0,G_\pi^J,G_1^J\}.
\end{equation}
The current lies in the previously established common pointwise and
Hilbert ball. Constants are independent of refinement and compatible
periodic volume. Here $D^*$ can be any specified physical fine
divergence, with its fixed orientation contractions. Taking
$\gamma=\beta+\varepsilon$, with fixed
$0<\varepsilon<1-\beta_0$, makes the constants uniform for
$0\le\beta\le\beta_0$. The source choice
$\varepsilon=(1-\beta_0)/2$ is included.
\end{theorem}

This is a bound on the right-differentiated operator, not an estimate
obtained by bounding $\norm{D^*\lambda}_\infty$. The constant can
grow as $\gamma\downarrow\beta$; the endpoint is not claimed.
The relevant source members are described at deposited Raw I,
lines 13601--13615 and 13759--13765.

\subsection{A remainder that accepts arbitrary bounded data}

Introduce the comparison operator
\begin{equation}\label{rd-massive-split}
 G_m=(\mathcal L+1)^{-1}I,\qquad R_m=G_0-G_m.
\end{equation}
It is an auxiliary decomposition, not a change of source precision.
Shrink the common complex chart so that $P_0+1$ is uniformly nonzero;
on the real chart it is at least one. The high denominators already
have this property. The high scalar diagonal of $R_m$ is
\begin{equation}\label{rd-high-remainder}
 \frac1{P_r}-\frac1{P_r+1}
    =\int_0^\infty(1-e^{-t})e^{-tP_r}\,dt,
                    \qquad r\ne0.
\end{equation}
The other terms are the finite-rank corrections of $G_0$ and the
bounded central subtraction $-(P_0+1)^{-1}\Pi_0$.

With one input derivative, one output derivative, and an output
$\beta$-H\"older increment, the small-time bound from the heat
estimates is
\[
 (1-e^{-t})t^{-(2+\beta)/2}\le t^{-\beta/2},
                       \qquad0<t\le1.
\]
It is integrable. Lower output members are easier, and large times
use the high-alias gap. In each rank term, the input derivative
multiplies its row by one fine multiplier. The high row
$\omega_s/P_s$ becomes $O(|\omega_s|/(1+|s|))$; the Woodbury row
is bounded by
$C(|a_s|/(1+|s|)^2+|\omega_s|/(1+|s|))$.
Their alias $\ell^2$ norms are uniform. The output
$C^{1,\beta}$ vectors retain the proved absolute alias sums for
$\beta<1$.

The full reconstructed remainder is periodic and holomorphic in
the resulting pointwise operator norms. Chart gluing and contour
shifting therefore give exponential coarse-cell kernels and
\begin{equation}\label{rd-remainder-bound}
 N_{\beta,\varpi}(R_mD^*f)
                  \le C_\beta\norm{f}_{\infty,\varpi},
                       \qquad0\le\beta<1.
\end{equation}
This remainder estimate genuinely accepts arbitrary bounded fine
data. It requires no H\"older extension across a cell boundary:
the right derivative has already been placed on its kernel.
Compatible periodic sampling gives the same bound on tori.

\subsection{Weighted heat moments down to the mesh scale}

Let $p_t^\eta(v)$ be the full-lattice probability heat kernel of
$\mathcal L$, with fine displacement $v\in\eta\mathbb Z^d$.
For $0<\gamma<1$ and $k=2,3$, we have
\begin{equation}\label{rd-heat-moment}
 \sum_v e^{\kappa|v|_1}|v|_1^\gamma
                   |D_\eta^kp_t^\eta(v)|
       \le Ct^{(\gamma-k)/2},\qquad0<t\le1.
\end{equation}
The probability-kernel convention includes the physical normalization:
if a kernel density is used instead, its $\eta^{-d}$ factor is
canceled by physical summation.

For $t\ge\eta^2$, apply the weighted heat variation estimate with
physical exponential scale $\delta=\kappa+t^{-1/2}$. Its lattice
parameter $\eta\delta$ lies in a fixed bounded interval. Semigroup
splitting into $k$ factors gives
\[
 \sum_v e^{\delta|v|_1}|D_\eta^kp_t^\eta(v)|
                \le C_kt^{-k/2},\qquad \eta^2\le t\le1,
\]
because $\delta^2t$ is bounded. The inequality
$|v|^\gamma\le C_\gamma t^{\gamma/2}e^{|v|/\sqrt t}$
proves \eqref{rd-heat-moment} in this range.

For $0<t<\eta^2$, expand the finite difference into its fixed list
of translates:
\[
 D_\eta^kp_t^\eta(v)
    =\eta^{-k}\sum_jc_jp_t^\eta(v+\eta a_j),
 \qquad |a_j|_1\le k,\quad\sum_j|c_j|\le C_k.
\]
The dimensionless random walk has time $t/\eta^2<1$, so its
exponential and fixed positive moments are uniformly bounded.
The weighted moment above is at most $C\eta^{\gamma-k}$.
Since $k>\gamma$ and $t<\eta^2$,
$\eta^{\gamma-k}\le t^{(\gamma-k)/2}$, proving
\eqref{rd-heat-moment} also below the mesh time.

For $t\ge1$, absorb $|v|^\gamma$ into a fixed additional
exponential weight $e^{\delta_1|v|}$. The weighted heat estimates
are bounded by $Ce^{c(\kappa+\delta_1)^2t}$. Choose
$\kappa,\delta_1$ small enough that this is at most $Ce^{t/2}$.
The massive factor $e^{-t}$ leaves an integrable tail. The required
zeroth and first heat derivatives obey the analogous estimates
without the polynomial moment.

\subsection{Cancellation on complete fields}

Every nonzero difference of the translation-invariant heat kernel
has exact zero row sum:
\begin{equation}\label{rd-zero-row}
 \sum_vD_\eta^kp_t^\eta(v)=0\qquad(k\ge1).
\end{equation}
The local weighted H\"older norm implies
\begin{equation}\label{rd-global-increment}
 \varpi(\operatorname{cell}x)|\lambda(y)-\lambda(x)|
 \le C\norm{\lambda}_{\gamma,\varpi}
        e^{\kappa|x-y|_1}|x-y|_1^\gamma.
\end{equation}
At bounded separation use a common fixed collar and neighboring
weight ratios. At larger separation use the two supremum bounds;
the polynomial factor is then bounded below and the cell-weight
ratio is controlled by physical distance up to a fixed offset.
No derivative of the weight is taken.

Subtract $\lambda(x)$ using \eqref{rd-zero-row}. Equations
\eqref{rd-heat-moment}--\eqref{rd-global-increment} give
\begin{align}
 \norm{D_\eta^2e^{-t\mathcal L}\lambda}_{\infty,\varpi}
   &\le Ct^{-1+\gamma/2}\norm{\lambda}_{\gamma,\varpi},
       \label{rd-heat-second}\\
 \norm{D_\eta^3e^{-t\mathcal L}\lambda}_{\infty,\varpi}
   &\le Ct^{-3/2+\gamma/2}\norm{\lambda}_{\gamma,\varpi},
                   \qquad0<t\le1.
       \label{rd-heat-third}
\end{align}
The difference between two second-derivative output values is
bounded by twice \eqref{rd-heat-second}, and by their physical
separation times \eqref{rd-heat-third}. Interpolating yields
\begin{equation}\label{rd-heat-holder}
 [D_\eta^2e^{-t\mathcal L}\lambda]_{\beta,\varpi}
       \le Ct^{-1+(\gamma-\beta)/2}
                                  \norm{\lambda}_{\gamma,\varpi}.
\end{equation}
Its integral is finite, with
\begin{equation}\label{rd-time-integral}
 \int_0^1t^{-1+(\gamma-\beta)/2}\,dt
                      =\frac2{\gamma-\beta}.
\end{equation}
For $\beta=0$ the supremum estimate suffices. The field member of
$G_mD^*\lambda$ uses only the first heat derivative, whose
$t^{-1/2}$ bound is integrable on bounded input. The gradient and
its H\"older member use
\eqref{rd-heat-second}--\eqref{rd-time-integral}. The large-time
massive tail is integrable as above. Therefore
\begin{equation}\label{rd-massive-bound}
 N_{\beta,\varpi}(G_mD^*\lambda)
                  \le C_{\beta,\gamma}
                                  \norm{\lambda}_{\gamma,\varpi}.
\end{equation}
All singular-integral cancellation has occurred in this global
translation-invariant calculation. On tori use the periodized heat
kernel, or periodically extend the fields and weights to the cover.
Periodic distance is no larger than cover distance, so the same
moment bounds and multiplicities give identical constants.

\subsection{Transfer to the actual current inverses}

Adding \eqref{rd-remainder-bound} and \eqref{rd-massive-bound}
proves \eqref{rd-main-estimate} for $G_0$. For either current
family, the proved weak-space perturbation obeys
\[
 N_{\beta,\varpi}(K_JA)
       \le C_*\norm{J}_\infty N_\varpi(A),\qquad
 C_*\norm{J}_\infty<\tfrac12.
\]
Apply the same convergent series to the baseline datum:
\begin{equation}\label{rd-current-transfer}
 G^JD^*\lambda=(I-K_J)^{-1}G_0D^*\lambda.
\end{equation}
The baseline $N$ bound requires positive input H\"older regularity;
the baseline $N_\beta$ bound uses $\gamma>\beta$. Every later term
is controlled by the established $N$-to-$N_\beta$ perturbation
bound. This proves the remaining assertions of
\eqref{rd-main-estimate}. The larger H\"older-data constant does
not shrink the current radius: it multiplies the initial datum,
not the contraction. Holomorphy follows from the same Neumann series.

Localization is legitimate with a smooth physical cutoff $\phi$
of fixed support collar and bounded first derivative. The product
difference identity gives
$\norm{\phi\lambda}_{\gamma,\varpi}
 \le C_\phi\norm{\lambda}_{\gamma,\varpi;\,	ext{collar}}$
without a power of $\eta^{-1}$. A source-centered exponential weight
then gives decay away from the support. A sharp cell indicator
would generally create $\eta^{-\gamma}$ H\"older jumps and is not
permitted by this argument. This matches the use of localized
cutoffs in the stated source member.

\subsection{The constrained-current member}

The source separately defines, at Raw I, line 13735,
\begin{equation}\label{rd-constrained-definition}
 \widetilde G^J=(I-\Gamma\Theta_J)^{-1}\Gamma,
 \qquad
 \Gamma=G_0-H_{\rm src}q_{\rm aux}^{-1}H_{\rm src}^*.
\end{equation}
The preceding section proved $Q\Gamma=0$ and, for a uniform
constant $B_\Gamma$,
$N_{\beta,\varpi}(\Gamma f)
 \le B_\Gamma\norm{f}_{\infty,\varpi}$.
It also proved
$\norm{\Theta_Jf}_{\infty,\varpi}
 \le\tau_\infty C_v\norm{J}_\infty
                      \norm{f}_{\infty,\varpi}$.
Decrease the common current radius to
\begin{equation}\label{rd-constrained-radius}
 \varepsilon_{\rm all}
  =\min\left\{\varepsilon_{\rm pt},
             \frac1{2B_\Gamma\tau_\infty C_v}\right\},
\end{equation}
omitting the second restriction if its coefficient is zero.
The series in \eqref{rd-constrained-definition} then converges in
weighted pointwise norms and gives the source (J6) estimate
\begin{equation}\label{rd-constrained-bound}
 N_{\beta,\varpi}(\widetilde G^Jf)
           \le2B_\Gamma\norm{f}_{\infty,\varpi},\qquad
                 Q\widetilde G^J=0.
\end{equation}
For the bound, estimate the outside $\Gamma$ in $N_\beta$ and
sum the remaining supremum-norm geometric series. For the exact
constraint, every term begins on the left with $\Gamma$, so applying
$Q$ gives zero term by term; the limit preserves the identity.
The map is analytic, and source-centered weights give exponential
coarse-cell locality. This statement keeps the source's constrained
inverse distinct from $G_\pi^J$ and $G_1^J$.

\subsection{Remaining distinctions}

The right-divergence theorem covers $G_0,G_\pi^J,G_1^J$ for the
stated H\"older input and also supplies the separate constrained
current member \eqref{rd-constrained-bound}. The deposited source
also names a right-divergence estimate for the distinct scalar
inverse $G'$. That member is not inferred by replacing $G'$ with
$G_0$. The following section gives its own scalar Woodbury adaptation
for a penalty in a fixed positive interval; verifying any separately
specified source penalty remains an application condition.

No current Laplacian estimate is asserted. The positive-loss argument
above will next be strengthened at each fixed equal exponent in
$(0,1)$. The results are homogeneous and flat,
with the corrected fixed-step constraint term already matched to
the deposited transcription. They do not prove nonflat source
stability, mixed retained-domain bounds, strict background locality,
or an interacting Yang--Mills ultraviolet limit.

% Independently internally audited strengthening of Notebook 12.
% Prefix ss-. Depends on rd heat estimates and flat alias notation.
\section{Equal-exponent Schauder bounds and the scalar inverse}
\label{ss-section}

The positive H\"older loss in the preceding right-divergence
theorem is sufficient for the displayed source estimates. It is
not necessary for a fixed exponent strictly between zero and one.
The following refinement is useful for nonlinear elliptic equations,
where an exponent loss at every iteration would not close a fixed
point argument.

\begin{theorem}\label{ss-equal-exponent}
For every fixed $0<\alpha<1$, the homogeneous flat operators obey
\begin{equation}\label{ss-main}
 N_{\alpha,\varpi}(GD^*\lambda)
       \le C_\alpha\norm{\lambda}_{\alpha,\varpi},
 \qquad G\in\{G_0,G_\pi^J,G_1^J\}.
\end{equation}
The spatial exponent is sufficiently small, and the current is in
the common ball for the $N$-to-$N_\alpha$ estimates established
above. Constants are uniform in refinement and compatible volume.
There is no uniform assertion as $\alpha$ tends to zero or one.
\end{theorem}

\begin{proof}
Use $G_0=G_m+R_m$ from \eqref{rd-massive-split}.
Equation \eqref{rd-remainder-bound} already controls $R_mD^*$
from bounded input into $N_\alpha$ for every $\alpha<1$.
Only the gradient H\"older member of $G_mD^*$ needs a new
argument. The heat moment and exact row cancellation, with
$\gamma=\alpha$, give
\[
 \norm{D^2e^{-t\mathcal L}\lambda}_{\infty,\varpi}
       \le Ct^{-1+\alpha/2}\norm{\lambda}_{\alpha,\varpi},
 \qquad
 \norm{D^3e^{-t\mathcal L}\lambda}_{\infty,\varpi}
       \le Ct^{-3/2+\alpha/2}\norm{\lambda}_{\alpha,\varpi}
\]
for $0<t\le1$, including times below the mesh scale.
For two fine output sites in a fixed collar at distance
$0<r\le1$, estimate their second-derivative difference by the
first bound for $t<r^2$, and by $r$ times the second for $t\ge r^2$.
The finite path used for this difference stays in a fixed enlarged
collar, so its weight comparison costs only a fixed constant.
Consequently the difference after integrating the massive heat
kernel is bounded, divided by the input norm, by
\begin{align}
 C\int_0^{r^2}t^{-1+\alpha/2}\,dt
 +Cr\int_{r^2}^1t^{-3/2+\alpha/2}\,dt+Cr
 &\le C\left(\frac2\alpha+
                     \frac2{1-\alpha}+1\right)r^\alpha.
 \label{ss-split-integral}
\end{align}
The last term is the large-time massive tail, using its third
derivative bound. Distances above one are controlled by the
supremum bound. The latter uses
$\int_0^1t^{-1+\alpha/2}dt=2/\alpha$ for the gradient and the
integrable first heat derivative for the field. This proves the
baseline assertion. It improves the preceding argument by
splitting the time integral for each output separation before
taking the H\"older supremum, rather than integrating a fixed
H\"older operator norm down to zero time.

The current transfer is exactly \eqref{rd-current-transfer}:
the baseline datum is in $N_\alpha$, and the perturbation maps
$N$ to $N_\alpha$ with a strict contraction in $N$. Its radius
does not depend on the initial datum constant. Periodic extension,
complete weighted fields and smooth localization work as before.
\end{proof}

\subsection{A separately normalized scalar statement}

Let
\[
 G'=(\mathcal L+Q_0^*a'Q_0)^{-1},\qquad
              0<a'_-\le a'\le a'_+<\infty,
\]
where the distinct scalar penalty ranges in a fixed compact
positive interval. This hypothesis is explicit: no unprinted
scale-dependent scalar penalty is assigned a value here.
With the scalar aliases $a_r$, put
\begin{equation}\label{ss-scalar-denominator}
 A=a_0a_0^-,\quad
 R_{s,1}=\sum_{r\ne0}\frac{a_ra_r^-}{P_r},\quad
 \widehat R_{s,1}=a'^{-1}+R_{s,1},\quad
 D_{s,1}=A+P_0\widehat R_{s,1}.
\end{equation}
The canceled scalar Woodbury formula is
\begin{align}
 G'_{00}&=\widehat R_{s,1}/D_{s,1},\nonumber\\
 G'_{0s}&=-a_0^-a_s/(P_sD_{s,1}),
 &G'_{r0}&=-a_r^-a_0/(P_rD_{s,1}),\nonumber\\
 G'_{rs}&=\frac{\delta_{rs}}{P_r}
       -\frac{P_0a_r^-a_s}{P_rP_sD_{s,1}},
 &&r,s\ne0.
 \label{ss-scalar-formula}
\end{align}
These formulas extend across $P_0=0$; no singular factor $1/P_0$
remains. They follow by the rank-one inverse identity away from
$P_0=0$ and analytic continuation. On the real compact chart,
$D_{s,1}\ge A\ge c>0$. Alias Parseval bounds and the high-alias
gap give uniform derivative bounds, so $D_{s,1}$ remains bounded
away from zero on a common complex strip. The constants may
depend on $a'_-,a'_+$.

Subtract the scalar massive inverse $G_m=(\mathcal L+1)^{-1}$.
The high diagonal has exactly the improved heat representation
\eqref{rd-high-remainder}. All other terms in
\eqref{ss-scalar-formula} have bounded central coefficients and
high output vector $a_r^-/P_r$. This vector has uniform absolute
alias summability after $1+\alpha$ physical differences when
$\alpha<1$: the bound
$|a_r|\le C\prod_\nu(1+|r_\nu|)^{-1}$ gives a dyadic shell sum
bounded by $C(1+n)^{d-1}2^{-n(1-\alpha)}$. One right derivative
leaves the input row bounded in alias $\ell^2$, since
$|u_ra_r/P_r|\le C|a_r|/(1+|r|)$ for high aliases.
The undifferentiated input rows have the same property.
Uniform cell Fourier conversion, periodic gluing and contour
shifting therefore give
\begin{equation}\label{ss-scalar-remainder}
 N_{\alpha,\varpi}\bigl((G'-G_m)D^*f\bigr)
       \le C_\alpha\norm{f}_{\infty,\varpi}.
\end{equation}
The bounded central massive subtraction is included in this
estimate. Only the complete reconstructed operator is periodized
and shifted; the individual central pieces need not be periodic.

Applying the same massive cancellation proof establishes
\begin{equation}\label{ss-scalar-right-divergence}
 N_{\alpha,\varpi}(G'D^*\lambda)
       \le C_\alpha\norm{\lambda}_{\alpha,\varpi},
                       \qquad0<\alpha<1.
\end{equation}
Using instead the preceding positive-loss proof also gives
$N_{\beta,\varpi}(G'D^*\lambda)
\le C_{\beta,\gamma}\norm{\lambda}_{\gamma,\varpi}$ for
$0\le\beta<\gamma<1$. This proves the scalar analytic estimate
for the stated penalty normalization. Application to a separately
specified source penalty requires verifying its uniform positive
interval. It neither replaces $G'$ by the one-form $G_0$ nor
identifies their gauge projections.

Equal H\"older exponents strictly between zero and one in
\eqref{ss-main} do not assert a Lipschitz-gradient endpoint or a
second-difference bound on arbitrary bounded forcing. The earlier
endpoint counterexample remains applicable to those other claims.

\begin{thebibliography}{9}
\bibitem{repo}
\href{https://github.com/ethanwsussman/Yang--Mills-LLM-experiment/tree/e392ef45013a581664bab76b90867170b30fe10e}
{Pinned Yang--Mills-LLM-experiment account}, commit
\texttt{e392ef45013a581664bab76b90867170b30fe10e}.
The homogeneous weak norms are Raw I:12740--12770; the bounded
current assumption is I:12951--12966; the distinct right-divergence
members are I:13601--13615; the current assertions and exclusion of
the current Laplacian member are I:13749--13804.
\bibitem{previous}
Companion research proofs: Notebook 09, \emph{The flat fine inverse
and a current neighborhood}; Notebook 11, \emph{The full flat
constraint-current family}. Source operator matching is to the pinned
deposited transcriptions. Original primary-paper arguments have not
all been independently reverified.
\end{thebibliography}
\end{document}
